% GENERATED FILE. Edit canonical lessons, then run npm run generate:books.
% canonical-source-hash: fnv1a64:22aca862bc0f54b1
% canonical-lessons: MW-C232-budho, MW-C232-mati, MW-C232-chikno, MW-C232-kakro, MW-C232-borri

\chapter{Old Earth, Smooth Stone}
\label{ch:mw-old-earth-smooth-stone}
\hlchaptermodality{232}

\begin{chapteropening}
\textbf{By the end of this chapter:} \emph{I can describe old people, the soil and the stones.}

Complete the last lesson of chapter 232: I can describe old people, the soil and the stones.
\end{chapteropening}

\section[\texorpdfstring{\mw{बूढो} (būḍho)}{बूढो (būḍho)}]{\mw{बूढो} (būḍho) — elderly; old (of a person)}

\label{lesson:MW-C232-budho}



\begin{quote}
Before the new one: say the Marwadi for round, then the Marwadi for loose.
\end{quote}

\subsection*{You'll want to know: \mw{बूढो}}
\textbf{\mw{बूढो}} — \emph{būḍho} — \textquotedblleft{}elderly; old (of a person)\textquotedblright{}.

\begin{quote}
\mw{दादा} \mw{बूढो} \mw{है।}
\end{quote}

\begin{quote}
\emph{dādā būḍho hai.}
\end{quote}

\begin{quote}
Grandfather is old.
\end{quote}

\begin{cousinweb}
is inherited from a reconstructed \emph{*buḍḍha}, \textquotedblleft{}old\textquotedblright{}; CDIAL lists Old Marwari \emph{būḍhaü}, \textquotedblleft{}old\textquotedblright{} (CDIAL 9271). \textbf{\mw{जूनो}} is old for things; \textbf{\mw{बूढो}} is old for people and animals.
\end{cousinweb}

\begin{grammarlens}[title={What you've built}]
One more word for this chapter.
\end{grammarlens}

\subsection*{Your turn}
\begin{itemize}
\raggedright
  \item \emph{Say it:} \emph{būḍho}
  \item \emph{Say it:} \emph{būḍho}, once more
  \item \emph{Say it:} \emph{dādā būḍho hai.} — Grandfather is old.
  \item \emph{Recall:} say the Marwadi for round, then the Marwadi for loose, then say \emph{būḍho} again
\end{itemize}

\begin{tcolorbox}[breakable,colback=teal!4,colframe=teal!35!black,title={Before you move on}]
What does \mw{बूढो} mean? (\textquotedblleft{}Elderly; old (of a person)\textquotedblright{}.) Say it once more.
\end{tcolorbox}
\section[\texorpdfstring{\mw{माटी} (māṭī)}{माटी (māṭī)}]{\mw{माटी} (māṭī) — soil; earth (the material)}

\label{lesson:MW-C232-mati}



\begin{quote}
Before the new one: say the Marwadi for loose, then the Marwadi for elderly.
\end{quote}

\subsection*{You'll want to know: \mw{माटी}}
\textbf{\mw{माटी}} — \emph{māṭī} — \textquotedblleft{}soil; earth (the material)\textquotedblright{}.

\begin{quote}
\mw{खेत} \mw{मांय} \mw{माटी} \mw{गीली} \mw{है।}
\end{quote}

\begin{quote}
\emph{khet mā̃y māṭī gīlī hai.}
\end{quote}

\begin{quote}
The soil in the field is wet.
\end{quote}

A feminine noun.

\begin{cousinweb}
\textbf{\mw{माटी}} (\emph{māṭī}) is inherited from Sanskrit \emph{mṛttikā}, \textquotedblleft{}earth, clay\textquotedblright{}; CDIAL lists Hindi and Gujarati \emph{māṭī}. CDIAL gives no Marwari form in this entry (CDIAL 10286).
\end{cousinweb}

\begin{grammarlens}[title={What you've built}]
One more word for this chapter.
\end{grammarlens}

\subsection*{Your turn}
\begin{itemize}
\raggedright
  \item \emph{Say it:} \emph{māṭī}
  \item \emph{Say it:} \emph{māṭī}, once more
  \item \emph{Say it:} \emph{khet mā̃y māṭī gīlī hai.} — The soil in the field is wet.
  \item \emph{Recall:} say the Marwadi for loose, then the Marwadi for elderly, then say \emph{māṭī} again
\end{itemize}

\begin{tcolorbox}[breakable,colback=teal!4,colframe=teal!35!black,title={Before you move on}]
What does \mw{माटी} mean? (\textquotedblleft{}Soil; earth (the material)\textquotedblright{}.) Say it once more.
\end{tcolorbox}
\section[\texorpdfstring{\mw{चिकणो} (chikṇo)}{चिकणो (chikṇo)}]{\mw{चिकणो} (chikṇo) — smooth; slippery}

\label{lesson:MW-C232-chikno}



\begin{quote}
Before the new one: say the Marwadi for elderly, then the Marwadi for soil.
\end{quote}

\subsection*{You'll want to know: \mw{चिकणो}}
\textbf{\mw{चिकणो}} — \emph{chikṇo} — \textquotedblleft{}smooth; slippery\textquotedblright{}.

\begin{quote}
\mw{भाटो} \mw{चिकणो} \mw{है।}
\end{quote}

\begin{quote}
\emph{bhāṭo chikṇo hai.}
\end{quote}

\begin{quote}
The stone is smooth.
\end{quote}

\begin{cousinweb}
\textbf{\mw{चिकणो}} (\emph{chikṇo}) is inherited from Sanskrit \emph{cikkaṇa}, \textquotedblleft{}slippery, unctuous\textquotedblright{}; CDIAL lists Hindi \emph{ciknā}, \textquotedblleft{}oily, smooth\textquotedblright{}, and a Gujarati cousin. CDIAL gives no Marwari form in this entry (CDIAL 4782).
\end{cousinweb}

\begin{grammarlens}[title={What you've built}]
One more word for this chapter.
\end{grammarlens}

\subsection*{Your turn}
\begin{itemize}
\raggedright
  \item \emph{Say it:} \emph{chikṇo}
  \item \emph{Say it:} \emph{chikṇo}, once more
  \item \emph{Say it:} \emph{bhāṭo chikṇo hai.} — The stone is smooth.
  \item \emph{Recall:} say the Marwadi for elderly, then the Marwadi for soil, then say \emph{chikṇo} again
\end{itemize}

\begin{tcolorbox}[breakable,colback=teal!4,colframe=teal!35!black,title={Before you move on}]
What does \mw{चिकणो} mean? (\textquotedblleft{}Smooth; slippery\textquotedblright{}.) Say it once more.
\end{tcolorbox}
\section[\texorpdfstring{\mw{कांकरो} (kā̃kro)}{कांकरो (kā̃kro)}]{\mw{कांकरो} (kā̃kro) — a pebble}

\label{lesson:MW-C232-kakro}



\begin{quote}
Before the new one: say the Marwadi for soil, then the Marwadi for smooth.
\end{quote}

\subsection*{You'll want to know: \mw{कांकरो}}
\textbf{\mw{कांकरो}} — \emph{kā̃kro} — \textquotedblleft{}a pebble\textquotedblright{}.

\begin{quote}
\mw{जूती} \mw{मांय} \mw{कांकरो} \mw{है।}
\end{quote}

\begin{quote}
\emph{jūtī mā̃y kā̃kro hai.}
\end{quote}

\begin{quote}
There's a pebble in my shoe.
\end{quote}

A masculine noun.

\begin{cousinweb}
\textbf{\mw{कांकरो}} (\emph{kā̃kro}) is inherited from Sanskrit \emph{karkara}, \textquotedblleft{}a stone\textquotedblright{}; CDIAL lists Hindi \emph{kā̃kar}, a nodule of limestone, and a Gujarati word for a pebble. CDIAL gives no Marwari form in this entry (CDIAL 2820).
\end{cousinweb}

\begin{grammarlens}[title={What you've built}]
One more word for this chapter.
\end{grammarlens}

\subsection*{Your turn}
\begin{itemize}
\raggedright
  \item \emph{Say it:} \emph{kā̃kro}
  \item \emph{Say it:} \emph{kā̃kro}, once more
  \item \emph{Say it:} \emph{jūtī mā̃y kā̃kro hai.} — There's a pebble in my shoe.
  \item \emph{Recall:} say the Marwadi for soil, then the Marwadi for smooth, then say \emph{kā̃kro} again
\end{itemize}

\begin{tcolorbox}[breakable,colback=teal!4,colframe=teal!35!black,title={Before you move on}]
What does \mw{कांकरो} mean? (\textquotedblleft{}A pebble\textquotedblright{}.) Say it once more.
\end{tcolorbox}
\section[\texorpdfstring{\mw{बोरड़ी} (borṛī)}{बोरड़ी (borṛī)}]{\mw{बोरड़ी} (borṛī) — a jujube bush}

\label{lesson:MW-C232-borri}



\begin{quote}
Before the new one: say the Marwadi for smooth, then the Marwadi for a pebble.
\end{quote}

\subsection*{You'll want to know: \mw{बोरड़ी}}
\textbf{\mw{बोरड़ी}} — \emph{borṛī} — \textquotedblleft{}a jujube bush\textquotedblright{}.

\begin{quote}
\mw{खेत} \mw{कनै} \mw{बोरड़ी} \mw{है।}
\end{quote}

\begin{quote}
\emph{khet kanai borṛī hai.}
\end{quote}

\begin{quote}
There's a jujube bush near the field.
\end{quote}

A feminine noun.

\begin{cousinweb}
\textbf{\mw{बोरड़ी}} (\emph{borṛī}) is inherited from Sanskrit \emph{badara}, \textquotedblleft{}the fruit of the jujube tree\textquotedblright{}; CDIAL lists Gujarati \emph{borṛī}, the jujube tree. CDIAL gives no Marwari form in this entry (CDIAL 9125).
\end{cousinweb}

\begin{grammarlens}[title={What you've built}]
One more word for this chapter.
\end{grammarlens}

\subsection*{Your turn}
\begin{itemize}
\raggedright
  \item \emph{Say it:} \emph{borṛī}
  \item \emph{Say it:} \emph{borṛī}, once more
  \item \emph{Say it:} \emph{khet kanai borṛī hai.} — There's a jujube bush near the field.
  \item \emph{Recall:} say the Marwadi for smooth, then the Marwadi for a pebble, then say \emph{borṛī} again
\end{itemize}

\begin{tcolorbox}[breakable,colback=teal!4,colframe=teal!35!black,title={Before you move on}]
What does \mw{बोरड़ी} mean? (\textquotedblleft{}A jujube bush\textquotedblright{}.) Say it once more.
\end{tcolorbox}
